% !PRINTED 394
\weberpage{394}
%<<CONT
of which the last two are conjugate imaginary, and which consequently,
since apart from $x, y, \xi, \xi'$ they have no common element, contain
only imaginary pairs besides, of which we have listed one, $\xi_1\eta_1$,
together with its conjugate $\xi_1'\eta_1'$.

It therefore suffices if we can prove the existence of one further real
double tangent. For this must then occur in 1), and must in this complex
belong to a real pair.

We now form in addition the two conjugate complexes:
\[
\begin{array}{rl}
4) & x\xi_1,\ \xi\eta_1 \\
5) & x\xi_1',\ \xi'\eta_1',
\end{array}
\]
in which $y$ certainly does not occur, since otherwise $x, \xi, y$ would be
azygetic (\S.\,96, 3.). For the same reason $\xi$ does not occur in 5),
nor $\xi'$ in 4).

Now if, \gesperrt{in the first case}, the complexes 4) and 5) are syzygetic, then
$\xi_1$ must occur in 5) and $\xi_1'$ in 4), and if 4) contains the pair $p\xi_1'$,
then the pair $p\xi_1$ must occur in 5). Thus $p$ is identical with its
conjugate element, that is, \gesperrt{real}.

If, \gesperrt{secondly}, the complexes 4) and 5) are azygetic, then $\xi_1, \xi$ do
not occur in 5), and $\eta_1$ must be contained in 5).
If $\eta_1$ is paired in 5) with an imaginary element $\xi_2'$, and
if in 4) yet another imaginary pair $\xi_3\eta_3$ occurs, then these
complexes continue in the following way:
\[
\begin{array}{rl}
4) & x\xi_1,\ \xi\eta_1,\ \xi_2\eta_1',\ \xi_3\eta_3,\ \xi_3'\eta_4 \\
5) & x\xi_1',\ \xi'\eta_1',\ \xi_2'\eta_1,\ \xi_3'\eta_3',\ \xi_3\eta_4',
\end{array}
\]
in which $\eta_4$ is also imaginary, and the accented letters are always
understood to mean the elements conjugate to the unaccented ones.
The two complexes then contain only one further pair each, $\xi_5\eta_5$,
$\xi_5'\eta_5'$, which must contain a common element. But it cannot
be that $\xi_5 = \eta_5'$, since otherwise also $\xi_5' = \eta_5$, and
hence the two pairs would be identical. Therefore it must be that $\xi_5 = \xi_5'$ (or
$\eta_5 = \eta_5'$); that is, one of these two double tangents
is \gesperrt{real}, and with this our theorem 2. is proved.

3. \gesperrt{If, besides the four real double tangents $x, y, p, q$,
whose existence theorem 2. asserts, yet another one is present,
then there exist eight real double tangents, which form four pairs
in a Steiner complex.}
